%% Three structural remarks on OpenAI's forced-blowup construction for Navier--Stokes (L0 scope) -- submission version (English)
%% Packaged from review\branch-D-note-skeleton.tex (frozen working draft, 893 lines, four independent
%% review rounds; targeted re-check r4 = CLEAR) and review\submission-note.tex (Chinese release draft).
%% This English version only translates prose, English-izes theorem/tag names, and re-typesets with
%% the standard article class. It changes NO mathematical assertion, NO line anchor, NO color-tag
%% structure, and NO L0/L1/L2 discipline.
%% Scope: L0 = a step of the device fails under its stated quantifiers, when the device is extended to
%%       the setting (including the external limit point j0=0, which lies outside the source's domain).
%%       L1 = no j0=0 variant within the class (H1)--(H3); valid only inside that class definition.
%%       L2 (existence/nonexistence of reflection-symmetric blowup) is NOT touched; such statements do not appear.
%% Line numbers: always refer to the reference text extraction navier-stokes.txt of the source PDF
%%       (10264 lines; SHA-256 given in the reproducibility paragraph). That extraction is a full-text
%%       copy of the source and is NOT redistributed with this note.
%% Color code: [Source] = verbatim in the extracted text; [Derived] = rewriting of source formulas;
%%       [Assumption] = an additional hypothesis or an unproven step.
%% Compile target: XeLaTeX + article (also verified with Tectonic 0.17.0 XeTeX engine).
%% This note is not peer-reviewed.
\documentclass[11pt,a4paper]{article}
\usepackage[margin=1in]{geometry}
\usepackage{amsmath,amssymb,amsthm,mathtools}
\usepackage{xcolor}
\usepackage{microtype}
\usepackage[hidelinks]{hyperref}
%% Do not hyphenate the source name.
\hyphenation{OpenAI}

%% Typesetting neatness: reduce occasional overfull boxes in math-heavy text and allow
%% long displays to break across pages (they affect only non-fatal warnings otherwise).
\setlength{\emergencystretch}{2em}
\allowdisplaybreaks
%% Bibliography set ragged-right: breakable URLs must not be stretched across a justified
%% line (otherwise inter-letter spacing inside a URL becomes visible).
\let\origthebibliography\thebibliography
\renewcommand{\thebibliography}[1]{\origthebibliography{#1}\raggedright}

\newtheorem{theorem}{Theorem}[section]
\newtheorem{lemma}[theorem]{Lemma}
\newtheorem{proposition}[theorem]{Proposition}
\newtheorem{corollary}[theorem]{Corollary}
\theoremstyle{definition}
\newtheorem{definition}[theorem]{Definition}
\newtheorem{claim}[theorem]{Claim}
\theoremstyle{remark}
\newtheorem{remark}[theorem]{Remark}

\newcommand{\R}{\mathbb{R}}
\newcommand{\supp}{\operatorname{supp}}
\newcommand{\lnos}[1]{(lines~#1)}
\newcommand{\ORIG}{\textcolor{blue}{\textbf{[Source]}}}
\newcommand{\DER}{\textcolor{teal}{\textbf{[Derived]}}}
\newcommand{\HYP}{\textcolor{orange}{\textbf{[Assumption]}}}
%% Uniform statement for steps that are only sketched:
\newcommand{\notproved}{\emph{This note does not prove this step; it is stated only as an observation
and does not constitute a complete proof.}}

\hypersetup{
  pdftitle={Three structural remarks on a forced finite-time blowup construction for Navier--Stokes},
  pdfauthor={Kaiyan Ren},
  pdfsubject={Three structural observations on a forced Navier--Stokes blowup construction},
  pdfkeywords={Navier--Stokes, forced finite-time blowup, five moments, terminal force, parity mechanism}
}

\title{Three structural remarks on a forced finite-time blowup construction for Navier--Stokes}
\author{Kaiyan Ren\thanks{Department of Computer Science and Engineering,
The Hong Kong University of Science and Technology, Hong Kong SAR, China.
\texttt{krenab@connect.ust.hk}}}
\date{September 10, 2026}

\begin{document}
\maketitle

\begin{abstract}
Under the standing assumption that the forced finite-time blowup construction of \mbox{OpenAI} for the
three-dimensional incompressible Navier--Stokes equations is correct as claimed, we record three
structural observations:
(i) in a fixed background the source's $O(N^{-1})$ bound for the five-moment increment is non-sharp,
the cumulative increment over the full modulation interval being $O(N^{-2})$;
(ii) the terminal force $F_0$ is rigidly equivalent to the truncated-shell commutator;
(iii) at $j_0=0$ a parity mechanism degenerates the endpoint cone test (B.19) of [S], and this
degeneration cannot be avoided within the class (analytic axis profile, vanishing leading residual stress
$T_0=0$, stress activated at $X_0$, the same cone test) (in-class rigidity; see Note 3 below).
Statements that a step fails are of level L0; the in-class rigidity holds only within its class
definition. We make no claim about the existence or nonexistence of reflection-symmetric forced blowup,
and we claim no new Navier--Stokes solution.
\end{abstract}

\tableofcontents

%% =====================================================================
\section{Introduction and scope}
\label{sec:intro}

\paragraph{The OpenAI construction (one paragraph).}
\ORIG
For every positive viscosity the paper constructs, from rest, a three-dimensional incompressible
Navier--Stokes solution with uniformly bounded kinetic energy and velocity unbounded in finite time,
driven by a smooth compactly supported force
\lnos{1--6,801--832}; the local field is given by Theorem~3.1 of [S]
(potential decomposition \eqref{eq:33}, flat residual \eqref{eq:34}, exterior heat profile \eqref{eq:35},
axis growth \eqref{eq:36})
\lnos{745--799};
after spatial--temporal truncation the field is localized to a compact set $K$, and the residual is
extended to a $C^\infty_c$ force (Prop.~10.1 and Lemma~10.2--10.3 of [S]) \lnos{7229--7386}.
The axis construction takes a small bias $0<j_0\le .05$ and $U^*=4\eta+j_0$
(\eqref{eq:B1}, \lnos{9002--9005});
the introduction presents this bias as a deliberate breaking of the $z\mapsto -z$ reflection symmetry
\lnos{184--190}.
The parameter order chooses $j_0$ before $\Lambda$ (\eqref{eq:B40}).

\paragraph{Scope of this note.}
\HYP
Throughout, we read the source formulas, quantifiers and parameter order under the hypothesis that the
construction is correct as claimed: we assume Theorem~3.1, the axis profile and endpoint test of
Appendix~B of [S], and the modulation--five-moment repair of \S4 Steps 2--3 hold as stated, and we then
make three observations about their internal structure.
This note does not re-prove the existence of the construction, does not connect the finite-dimensional
moment maps to the full fluid field, and claims no new Navier--Stokes solution or regularity criterion.

\paragraph{Scope discipline (L0).}
\DER
Outward statements use level L0 only:
\emph{when the device is extended to the stated setting (including the external limit point $j_0=0$,
which lies outside the source's domain, the source allowing only $0<j_0\le.05$ \lnos{9002}),
a step fails under the device's own quantifiers.}
The statement that there is no $j_0=0$ variant within one function class (L1) has been recorded as
Prop.~\ref{prop:L1} (the class definition (H1)--(H3) is part of its hypotheses);
the existence or nonexistence of reflection-symmetric forced blowup (L2) is not touched by this note,
and no such statement appears below.
The first cone inequality $a-b_s w>0$ ($w=N_s/(EQ_s)$\lnos{8323}) and the second cone test $v_s>2$
are kept separate and are not conflated into one ``cone''.

\paragraph{Color code and line references.}
\ORIG is verbatim readable in \texttt{navier-stokes.txt};
\DER is a rewriting of source formulas under the stated setting;
\HYP is an additional hypothesis or an unproven step.
Every assertion carries line numbers. The internal working record used during the writing (not
distributed) contains no independent five-moment $M$ formula or closed form of $F_0$;
the formulas of Notes 1 and 2 are rewritten from the cited source lines.
Paragraphs of that record marked as superseded (SUPERSEDED) are not cited at all.

\paragraph{Notation, line references and reproducibility.}
Every \lnos{xxx} refers to line numbers of the reference text extraction \nolinkurl{navier-stokes.txt}
of the source PDF (\cite{source}): a plain-text extraction with 10264 lines, 507\,517 bytes, each page
separated by a marker of the form \nolinkurl{=== PAGE n ===} (165 page markers in all) and no ASCII
form-feed characters, with SHA-256
\begin{center}
\texttt{76a236df4f12ab2261f07047ae56a1dcc320b515c55c82653fea48942b10d674.}
\end{center}
The extraction was based on \texttt{pdftotext}; the version is unknown (the existing extraction carries
no version metadata), and the page-marker format is not the \texttt{pdftotext} default (which paginates
with form feeds), so a fresh \texttt{pdftotext} run will not reproduce this file byte for byte.
Being a full-text copy of \cite{source}, the extraction is not redistributed with this note; the
fingerprint above identifies it unambiguously, and the author will supply it on request.
A reader working from the source PDF directly can still locate almost every \ORIG statement without it,
because this note prints the source's own tags next to the line anchors: equation numbers such as
(B.19) or (10.9), statement numbers such as Prop.~B.3 of [S], and, for the extraction, the page markers.

\paragraph{Numbering and citation convention.}
The theorem-like environments of this note (Theorem/Lemma/Proposition/Corollary/Definition/Claim/Remark)
are numbered consecutively by ``section.counter'' and displayed formulas by $(1),(2),\dots$.
All numbers that refer to the source paper are written with ``of [S]''
(Theorem~3.1 of [S], Prop.~B.3 of [S], Lemma~B.4 of [S], Cor.~B.6 of [S], Theorem~4.6(iii) of [S], etc.)
to avoid confusion with the numbering of this note. Parenthesized numbers carrying an appendix letter
or a section/paragraph, such as (B.19), (B.26), (10.9), (A.24), always refer to the source formulas;
where this note re-typesets a source formula it uses its own number and states the source's tag and
line numbers immediately adjacent.

\paragraph{Status of proofs.}
This note records three structural observations; it is not a research paper giving complete proofs.
Prop.~\ref{prop:osc}--\ref{prop:V}, \ref{prop:shell}--\ref{prop:rigid},
\ref{prop:p2zero}--\ref{prop:B19deg} give only argument sketches: every proof block headed
``Argument sketch'' ends with the explicit statement ``This note does not prove this step; it is stated
only as an observation and does not constitute a complete proof.''
Propositions without a proof block (\ref{prop:rigid}, \ref{prop:p1}, \ref{prop:B19deg}) have their
derivation in the steps listed in their statements, likewise not constituting a complete proof.
The argument of Prop.~\ref{prop:L1} is its items 1--3, which depend on the sketch-only propositions
above, and is valid only within the class definition (H1)--(H3).
Steps tagged \HYP are unproven and must not be used as established conclusions;
the evidence level is read off the color code: \ORIG is checkable verbatim against the extracted text,
\DER is a rewriting of source formulas, \HYP is unproven.

%% ---------------------------------------------------------------------
\subsection*{Notation: source coordinates and symbols}
\label{sec:notation}

\ORIG
Similarity coordinates and weights \lnos{1343--1350}:
\begin{equation}
\label{eq:41}
\begin{aligned}
\tau&=1-t,\quad
A=\tfrac12+h,\quad
D=\tfrac12-h,\quad
0<h<\tfrac12,\\
z&=q^D\eta,\quad
\tau=q(1-\eta^2),\quad
d=1-\eta^2,\quad
L=1-2h\eta^2,\quad
s=r^2/2,\quad
X=s/q.
\end{aligned}
\end{equation}
Hence at $\eta=0$ one has $d=L=1$.
Profiles \lnos{1374--1378}:
\begin{equation}
\label{eq:43}
u^{(0)}_\theta=q^{-A}E,\quad
u^{(0)}_z=q^{-A}U,\quad
r u^{(0)}_r=V_0,\quad
p^{(0)}=q^{-2A}\Pi,\quad
E=C^{-1}\sqrt{2X}\,\varphi.
\end{equation}
Stress coefficients \lnos{1488--1496}:
\begin{equation}
\label{eq:411}
p_s=\Bigl(\frac{XQ_s}{L},\frac{XN_s}{LE}\Bigr),\quad
s=(a,-b_s),\quad
a=1-2D_X\log E,\quad
b_s=\frac{2D_X U}{E}.
\end{equation}
First cone direction $w=N_s/(EQ_s)$\lnos{8323}; the first cone inequality $a-b_s w>0$
(the sufficient test of Lemma~4.5 of [S] and (A.24)\lnos{1778,8324--8325}).
Below, $p_1=p_{s,1}$, $p_2=p_{s,2}$. The shear $v_s=a(1+t_s^2)=a+b_s^2/a$\lnos{1748--1751};
in the axis region and the reference profile, where $T_0=0$ (i.e.\ $p_s=s$),
$v_s=p_1+p_2^2/p_1$ (for $p_1>0$; the $p$-form appears in \lnos{9385--9389}).

%% =====================================================================
\section{Note 1: non-sharpness of the \texorpdfstring{$O(N^{-1})$}{O(N\textasciicircum -1)} five-moment bound}
\label{sec:note1}

\subsection{Source statements}
\label{sec:note1-orig}

\ORIG
The zero-mean primitives $A,B$ of the periodic shear satisfy \lnos{2390--2401}
\begin{equation}
\label{eq:AB}
\partial_\varphi A=-\tfrac12(a_L-a_0),\qquad
\partial_\varphi B=\tfrac12 E_0(b_L-b_0),\qquad
\int_0^1 A\,d\varphi=\int_0^1 B\,d\varphi=0,
\end{equation}
and vanish near the two ends of $I=[X_-,X_+]$, where the loop is constant, with zero extension off $I$.
The modulation \lnos{2402--2408}:
\begin{equation}
\label{eq:438}
E_N=E_0\exp\Bigl(\frac{A(X,\eta,N\log X)}{N}\Bigr),\qquad
U_N=U_0+\frac{B(X,\eta,N\log X)}{N}.
\end{equation}
Pointwise and shear bounds \lnos{2421--2426}:
\begin{equation}
\label{eq:439}
\|U_N-U_0,\,E_N-E_0\|_k\le A_k/N,\qquad
\|a_N-a_L,\,b_N-b_L\|_k\le B_k/N.
\end{equation}
The exact integral of the five-moment increment (fixed interval away from zero) \lnos{2428--2448}:
put $u_N=U_N-U_0$, $e_N=E_N-E_0$, $\Delta m_N=m_N-m_0$,
\begin{equation}
\label{eq:Dm}
\Delta m_N(X,\eta)
=\int_{X_-}^{X}
\begin{pmatrix}
u_N \\
\sqrt{2x}\,e_N \\
H_0 u_N+U_0\sqrt{2x}\,e_N+\sqrt{2x}\,u_N e_N \\
2U_0 u_N+u_N^2-E_0 e_N-e_N^2/2 \\
E_0 e_N/x+e_N^2/(2x)
\end{pmatrix}
dx.
\end{equation}
The product rule and \eqref{eq:439} give \lnos{2449--2452}
\begin{equation}
\label{eq:Dm-orig}
\|\Delta m_N\|_k\le D_k/N,\qquad
\|\Pi_N-\Pi_0^{\mathrm{prof}}\|_k=\|\Delta C_{p,N}\|_k\le D_k/N.
\end{equation}
Then take $d_N=-\Delta m_N(Y_1,\eta)$, arranged as $(M,J,I,S,C_p)$ \lnos{2497--2511}:
\begin{equation}
\label{eq:442}
B(\eta)c+Q_\eta(c,c)=d_N(\eta).
\end{equation}
and require $N\ge 8\max_{k=0,2}\nu_k^2 q_k D_k$
so that $\|c\|_{C^2}\le 2\nu_2 D_2/N$\lnos{2546--2552}.

\begin{remark}[Ordering of rows]
\ORIG
The five integrands of \eqref{eq:Dm} are arranged as $(M,I,J,S,C_p)$;
the target vector of \eqref{eq:442} is arranged as $(M,J,I,S,C_p)$\lnos{2497--2511}.
The two orderings must not be mixed when citing.
\end{remark}

\subsection{Non-sharpness: cumulative moments improve to \texorpdfstring{$O(N^{-2})$}{O(N\textasciicircum -2)}}
\label{sec:note1-der}

\begin{proposition}[Oscillatory-integral boundary term of order $O(N^{-1})$]
\label{prop:osc}
\DER
Let $[a,b]\subset(0,\infty)$ be fixed, let $P(X,\eta,\varphi)$ be smooth, $1$-periodic in $\varphi$ with
zero mean, and let $Q$ be its periodic primitive ($\partial_\varphi Q=P$). The phase $\varphi=N\log X$
is independent of $\eta$ \lnos{2409--2410,2421}. Then
\begin{equation}
\label{eq:ibp}
P(X,\eta,N\log X)
=\frac{X}{N}\Bigl(\partial_X Q(X,\eta,N\log X)-Q_X(X,\eta,N\log X)\Bigr),
\end{equation}
so that for a smooth weight $g$ on a fixed interval,
\[
\int_a^x g P_N\,dX
=\frac{[Xg Q_N]_a^x}{N}
-\int_a^x\frac{(Xg)'Q_N+Xg(Q_X)_N}{N}\,dX
\]
is $O(N^{-1})$ uniformly in $x$. The same holds for any fixed finite order of $\eta$-derivatives
(constants may depend on the fixed background, the interval and the order of differentiation).
\end{proposition}

\begin{proof}[Argument sketch]
The chain rule gives \eqref{eq:ibp}; the integration by parts is standard.
The complete $C^k_\eta$ estimate (only derivative bounds on a compact parameter set are needed) is not
written out in this note.
\notproved
\end{proof}

\begin{proposition}[Cumulative five moments $O(N^{-2})$]
\label{prop:N2}
\DER
In a fixed background $(U_0,E_0)$ and fixed $k$, put $u_N=B_N/N$ and
$e_N=E_0\bigl(\exp(A_N/N)-1\bigr)$. From \eqref{eq:AB} and \eqref{eq:438},
\begin{equation}
\label{eq:e-exp}
e_N=E_0\frac{A_N}{N}+O_{C^k_\eta}(N^{-2}).
\end{equation}
All linear increments in \eqref{eq:Dm} are $N^{-1}$ times zero-mean oscillatory integrals;
the remaining terms are quadratic or higher. Hence uniformly on $J=[X_-,Y_1]$,
\begin{equation}
\label{eq:N2}
\sup_X\|\Delta m_N(X,\cdot)\|_{C^k_\eta}\le C_k N^{-2}.
\end{equation}
Thus the source's $O(N^{-1})$ bound \eqref{eq:Dm-orig} is non-sharp for the \emph{cumulative} five moments.
\end{proposition}

\begin{proof}[Argument sketch]
Substitute \eqref{eq:e-exp} into \eqref{eq:Dm}. Linear terms use Prop.~\ref{prop:osc};
quadratic terms are already $O(N^{-2})$. The functions $A,B$ vanish near the ends of $I$\lnos{2401},
so the boundary terms are compatible with the choice of interval.
The explicit componentwise constants $C_k$ are not given in this note.
\notproved
\end{proof}

\begin{proposition}[Second-order leading term $V$ on the full interval]
\label{prop:V}
\DER
Because the oscillatory functions vanish near the ends, a smooth zero-mean linear term can be integrated
by parts repeatedly, and the full-interval integral decays super-algebraically at any fixed order.
Let $\langle\,\cdot\,\rangle$ denote the $\varphi$-average and integrate over $I$. Arranged as
$(M,I,J,S,C_p)$, in $C^k_\eta$ for a fixed background,
\[
N^2\Delta m_N(Y_1,\eta)\to V(\eta),
\]
\begin{equation}
\label{eq:V}
\begin{aligned}
V_M&=0,\\
V_I&=\int\sqrt{2X}\,E_0\,\frac{\langle A^2\rangle}{2}\,dX,\\
V_J&=\int\sqrt{2X}\,E_0\,U_0\frac{\langle A^2\rangle}{2}\,dX\\
&\qquad+\int\sqrt{2X}\,E_0\,\langle AB\rangle\,dX,\\
V_S&=\int\bigl(\langle B^2\rangle-E_0^2\langle A^2\rangle\bigr)\,dX,\\
V_{C_p}&=\int\frac{E_0^2\langle A^2\rangle}{X}\,dX.
\end{aligned}
\end{equation}
In particular: if $A\not\equiv 0$ for some $\eta$, then $V_{C_p}>0$;
if $A\equiv 0$ but $B\not\equiv 0$, then $V_S=\int\langle B^2\rangle>0$.
For a nontrivial modulation the five-moment vector generally has a genuine second-order mean error,
so that $O(N^{-2})$ is sharp as a vector bound for generic modulations.
\end{proposition}

\begin{proof}[Argument sketch]
Expand the exponential to second order; the cubic remainder integrates to $O(N^{-3})$;
the quadratic terms split into a periodic mean and a zero-mean part, the latter integrated by parts again.
\emph{Only on the full interval can the linear terms be made super-algebraically small; a partial
cumulative (moving endpoint) carries boundary terms and cannot use $V$ as is.}
\notproved
\end{proof}

\subsection{What the improvement does not cover}
\label{sec:note1-limit}

\begin{claim}
\label{cl:not-pointwise}
\DER
Prop.~\ref{prop:N2} improves the integrated quantity, not the pointwise profile bound $O(N^{-1})$
or the shear error $O(N^{-1})$ of \eqref{eq:439}. Subsequent changes of $p$ may still contain the
contribution of the profile itself and cannot all be replaced by $O(N^{-2})$.
\end{claim}

\begin{claim}
\label{cl:not-uniform}
\HYP
If the normalization constants are also controlled, one may write $|d_2-d_1|\le C_k(\lambda)N^{-2}$.
For fixed $\lambda$ and $N$ large enough this is indeed smaller than $\lambda/N$.
This is not a two-parameter uniform statement. From this upper bound one may not conclude that
``only $\lambda N\gtrsim 1$ suffices'', nor take the scaling $N\sim\lambda^{-2}$ suggested by
$\nu_k=O(\lambda^{-1})$ as an established asymptotics of the full construction
(the source fixes $\lambda$ before increasing $N$\lnos{2323--2326,2542--2545}).
\notproved
\end{claim}

\begin{remark}[The $\lambda$-factor in a normal combination: observation only]
\HYP
For fixed constants $G,c_0$, the second-order leading term of the combination
$\Delta J-G\Delta I-c_0\Delta M$ is
\[
\int\sqrt{2X}\,E_0\Bigl((U_0-G)\frac{\langle A^2\rangle}{2}+\langle AB\rangle\Bigr)\,dX.
\]
Whether this expression contains a structural $\lambda$-factor is not proven. The source's repair target
is $-\Delta m$, of opposite sign.
\notproved
\end{remark}

%% =====================================================================
\section{Note 2: the terminal shell force \texorpdfstring{$F_0$}{F\_0}, explicit formula and rigid equivalence}
\label{sec:note2}

\subsection{Localization and the definition of the force}
\label{sec:note2-loc}

\ORIG
Potential decomposition and exterior heat profile of Theorem~3.1 of [S] \lnos{745--790}:
\begin{equation}
\label{eq:33}
u=\operatorname{curl} A+B e_\theta,
\end{equation}
\begin{equation}
\label{eq:34}
\bigl|\partial_x^\alpha\partial_t^b R(u,p)\bigr|
\le C_{\alpha,b,N,X_1} q^N
\qquad(0\le X\le X_1,\;q\downarrow 0),
\end{equation}
\begin{equation}
\label{eq:35}
X\ge X_{\mathrm{ext}}\ \Rightarrow\
A=0,\quad
B=K(r,\tau),\quad
p=-\int_r^\infty\frac{K(\rho,\tau)^2}{\rho}\,d\rho,\quad
K(r,\tau)=r^{-1-2h}H_{\mathrm{ext}}(\tau/r^2).
\end{equation}
Axis growth \lnos{791--799}:
\begin{equation}
\label{eq:36}
u_\theta\bigl(\sqrt{2X_{\mathrm{in}}\tau},\,0,\,0,\,1-\tau\bigr)
=\tau^{-A}\bigl(e_0+O(\tau^{2h})\bigr)\qquad(\tau\downarrow 0).
\end{equation}

\ORIG
Localization of \S3.5 of [S] \lnos{801--832}: take $c=\chi_x\chi_t$,
with $\chi_x$ axisymmetric and supported in the compact set $K$, equal to $1$ near $(0,1)$,
and $\chi_t$ vanishing on $[0,1-\tau_0]$ and equal to $1$ where $1-t\le\tau_0/2$.
\begin{equation}
\label{eq:loc}
u=\operatorname{curl}(cA)+c B e_\theta
=c\,u_{\mathrm{loc}}+\nabla c\times A,\qquad
p=c\,p_{\mathrm{loc}}.
\end{equation}
For $t<1$ put $f=R(u,p)$. The difference $f-c R(u_{\mathrm{loc}},p_{\mathrm{loc}})$ contains
$(\partial_t c-\Delta c)u_{\mathrm{loc}}$, $p_{\mathrm{loc}}\nabla c$,
derivatives and advection of $\nabla c\times A$, and $(c^2-c)(u_{\mathrm{loc}}\cdot\nabla)u_{\mathrm{loc}}$,
supported in the truncation transition region, away from a neighborhood of $(0,1)$\lnos{815--819}.

\ORIG
Prop.~10.1 of [S] \lnos{7229--7258}: there is a compact set $K\subset\R^3$ such that
$\operatorname{div} u=0$, $u(\cdot,t)=p(\cdot,t)=0$ for sufficiently small $t\ge 0$, and
$(u,p)$ agrees with $(u_{\mathrm{loc}},p_{\mathrm{loc}})$ in a spatial neighborhood of the origin for $t$
sufficiently close to $1$. Definition of the force \lnos{7259--7260}:
\begin{equation}
\label{eq:105}
f=R(u,p)=\partial_t u+(u\cdot\nabla)u-\Delta u+\nabla p\qquad(0\le t<1).
\end{equation}

\subsection{The terminal value \texorpdfstring{$F_0$}{F\_0} and the shell formula}
\label{sec:note2-formula}

\ORIG
Lemma~10.2 of [S] \lnos{7275--7287}: for every spatial multi-index $\alpha$ and integer $j\ge 0$,
$\partial_x^\alpha\partial_t^j f$ converges uniformly on $\R^3$ as $t\uparrow 1$.
There exist $F_j\in C^\infty_c(\R^3;\R^3)$, all supported in $K$, such that
\begin{equation}
\label{eq:106}
\lim_{t\uparrow 1}\partial_x^\alpha\partial_t^j f(x,t)
=\partial_x^\alpha F_j(x),\qquad
\partial_x^\alpha F_j(0)=0.
\end{equation}
Below we write
\begin{equation}
\label{eq:F0}
F_0:=F_{j=0}=\lim_{t\uparrow 1}f(\cdot,t)\in C^\infty_c(K;\R^3).
\end{equation}

\ORIG
Exterior heat profile \lnos{7312--7324}: where $r$ is bounded below and $q$ is sufficiently small
(so that $X\ge X_{\mathrm{ext}}$),
\begin{equation}
\label{eq:107}
K(r,\tau)=c_\infty s^{-A}H(2\tau/s),\qquad s=r^2/2,
\end{equation}
with $H_{\mathrm{ext}}(Z)=2^A c_\infty H(4Z)$. For every $j$,
\begin{equation}
\label{eq:108}
\bigl|\partial_\tau^j K(r,\tau)\bigr|\le C_j r^{-1-2h-2j},\qquad
\Bigl|\partial_\tau^j\frac{K(\rho,\tau)^2}{\rho}\Bigr|\le C_j\rho^{-3-4h-2j}.
\end{equation}
The untruncated exterior residual vanishes identically\lnos{7343--7344}.

\ORIG
Near the origin the truncation equals $1$. We use \eqref{eq:34} on $X\le X_{\mathrm{ext}}$,
and the residual vanishes identically for larger $X$; so for every $\alpha,j,N$, on a fixed
neighborhood of $(0,1)$ \lnos{7349--7356}
\begin{equation}
\label{eq:109}
\bigl|\partial_x^\alpha\partial_t^j f(x,t)\bigr|
\le C_{\alpha,j,N}\bigl(\tau+|z|^{1/D}\bigr)^N.
\end{equation}
Hence $F_0$ vanishes to infinite order at the origin.

\begin{proposition}[Shell-commutator representation]
\label{prop:shell}
\DER
Take $t$ sufficiently close to $1$ so that $\chi_t\equiv 1$ (i.e.\ $1-t\le\tau_0/2$\lnos{804,7248--7249});
then $c=\chi_x$. Hence
\begin{equation}
\label{eq:F0-c}
F_0(x)
=\lim_{t\uparrow 1}
R\bigl(\operatorname{curl}(\chi_x A)+\chi_x B e_\theta,\,\chi_x p_{\mathrm{loc}}\bigr).
\end{equation}
Split by support:
\begin{enumerate}
\item \emph{The kernel} (the neighborhood where $\chi_x\equiv 1$): $F_0$ is the terminal value of the
local residual, flat (zero to infinite order) at $0$ by \eqref{eq:109}.
\item \emph{The untruncated exterior} ($X\ge X_{\mathrm{ext}}$ with $\chi_x\equiv 1$):
the untruncated residual vanishes identically\lnos{7343--7344}, so $F_0=0$.
\item \emph{The truncation shell} $\operatorname{supp}|D\chi_x|$
(which may meet the time-truncation transition where $q$ is bounded below\lnos{821--825}):
$F_0$ equals the terminal value of $f-c R(u_{\mathrm{loc}},p_{\mathrm{loc}})$, i.e.\ the commutator
listed in [S]\lnos{815--818}.
On the exterior shell, where $A=0$, this commutator involves only the terminal heat flow $K(\cdot,0)$,
the terminal pressure integral and the fixed cutoff $\chi_x$; its $\tau$-derivative bounds are controlled
by \eqref{eq:108}.
\end{enumerate}
\end{proposition}

\begin{proof}[Argument sketch]
For $\chi_t\equiv 1$, \eqref{eq:loc} and \eqref{eq:105} give \eqref{eq:F0-c}.
The three-way split follows from the definition of $\chi_x$\lnos{7246--7247}, the flatness
\eqref{eq:109} and the vanishing exterior residual \lnos{7343--7344}.
The componentwise Cartesian expansion of the commutator (including the vanishing of
$\nabla\chi_x\times A$ when $A=0$) is not written out in this note.
\notproved
\end{proof}

\begin{remark}[Time extension is not the shortest transient]
\ORIG
Lemma~10.3 of [S] \lnos{7384--7395} uses an order-shrinking time truncation to realize
\begin{equation}
\label{eq:1011}
f(x,1+\sigma)=\sum_{j=0}^\infty\chi_0(b_j\sigma)\frac{\sigma^j}{j!}F_j(x),
\end{equation}
supported in $K\times[0,2]$.
\DER
This gives the existence of a smooth compactly supported force, not the shortest switching transient.
\HYP
Whether shrinking the $\chi_t$ transition width breaks the derivative compatibility of \eqref{eq:106}
currently has no anchor.
\notproved
\end{remark}

\subsection{Rigid equivalence}
\label{sec:note2-rigid}

\begin{proposition}[Rigid equivalence with the shell data]
\label{prop:rigid}
\DER
Under the truncation $c=\chi_x\chi_t$ and the local field as in Theorem~3.1 of [S], the following
three data determine one another (rigid equivalence):
\begin{enumerate}
\item the terminal force $F_0=\lim_{t\uparrow 1}f$;
\item the full family of Taylor coefficients $\{F_j\}_{j\ge 0}$ of Lemma~10.2 of [S]
(via the time compatibility \eqref{eq:1010});
\item the terminal value of the commutator on the truncation shell
(flat-zero kernel $+$ zero untruncated exterior residual $+$ shell commutator).
\end{enumerate}
The equivalence does not depend on non-flat velocity details of the axis kernel: \eqref{eq:109}
suppresses derivatives of every order near the origin.
\end{proposition}

\ORIG
Time compatibility \lnos{7367--7378}:
\begin{equation}
\label{eq:1010}
\partial_x^\alpha\partial_t^j f(x,t)
=\partial_x^\alpha F_j(x)
-\int_t^1\partial_x^\alpha\partial_t^{j+1}f(x,s)\,ds.
\end{equation}

\begin{claim}[Whether $F_0$ encodes $j_0$: the exterior factor carries no bias; the fingerprint question is open]
\label{cl:F0-j0}
\DER
The bias $j_0$ appears in the axis velocity $U^*=4\eta+j_0$\lnos{9003--9004},
and does not appear in the exterior heat-flow factor of the shell commutator \eqref{eq:F0-c}.
Flatness \eqref{eq:109} suppresses the contribution of the kernel velocity to $F_0$ at the origin to
infinite order\lnos{7349--7356}.
\HYP
``$F_0$ cannot by itself serve as a fingerprint of a symmetric/asymmetric variant'' \emph{does not}
follow from the above facts: (10.9) suppresses only a neighborhood of the origin; the kernel region
$X<X_{\mathrm{ext}}$ inside $K=\operatorname{supp}\chi_x$\lnos{7244--7247} has positive measure, and the
source states only that the residual vanishes identically for $X\ge X_{\mathrm{ext}}$\lnos{779}; the
dependence of the terminal residual in that kernel on $j_0$ is not addressed, so the fingerprint question
is open.
To distinguish variants by the terminal field one would need additional normalized observables on
$C_\tau$ (e.g.\ the axis value $U(0,0)$, the $e_0$ of \eqref{eq:36}, or $E(X_*,0)$ at $\eta=0$), and a
proof that the high-frequency corrections of \S5 and \S9 of [S] do not erase that observable. This step
is not done.
\notproved
\end{claim}

\begin{remark}[Interior points in a bounded domain]
\ORIG
Placing the localized $F_0$ in a bounded smooth domain with $0$ in the interior is the trivial
restriction of Prop.~10.1 of [S]\lnos{7229--7232}: the domain need only contain $K$.
This note does not regard that restriction as a nontrivial boundary blowup.
\end{remark}

%% =====================================================================
\section{Note 3: the load-bearing role of the bias \texorpdfstring{$j_0$}{j\_0} --- a parity mechanism}
\label{sec:note3}

\subsection{Setup and the source endpoint test}
\label{sec:note3-setup}

This section fixes \HYP (the same test point as the second review round):
$j_0=0$, $U^*=4\eta$, endpoint $(X,\eta)=(4/\Lambda,0)$, $Y=\Lambda X=4$.

\ORIG
Axis data \lnos{8998--9009}:
\begin{equation}
\label{eq:B1}
U^*=4\eta+j_0,\qquad
H^*=D\eta+d U^*,\qquad
W^*=1-d U^*_\eta-2D\eta U^*,
\end{equation}
with
\[
-W^*=3-8h\eta^2+(1-2h)j_0\eta.
\]
\DER
Substituting $j_0=0$, $\eta=0$ into \eqref{eq:B1} gives $W^*(0)=-3$, $H^*(0)=0$, hence
$\eta_0=0$\lnos{9004--9008}.
\ORIG
$\Pi_0$ is even\lnos{8998--9000}. The dichotomy \eqref{eq:B2} \lnos{9014--9019}:
\begin{equation}
\label{eq:B2}
\begin{gathered}
\text{first choose }\delta_*>0\text{ so that }\{|Z^*|\le\delta_*\}\text{ avoids a neighborhood of }\eta_0,\\
\text{then choose }\sigma_*\text{ so that }\chi>.99\text{ on }\{|Z^*|\le\delta_*\}.
\end{gathered}
\end{equation}
The amplitude $C$ is chosen after $\Lambda$\lnos{9037--9038,9215--9218};
Lemma~B.4 of [S] ``first $\Lambda$, then $C$, sufficiently large''\lnos{9383};
\S B.7 of [S]: increasing $C$ moves $X_R$ outward without lengthening the transition\lnos{9580--9586}.
The absence of an upper bound is thus pinned.
Azimuthal data \lnos{9027--9036}:
\begin{equation}
\label{eq:B3}
\zeta^*=-\frac{L H^*}{H^{*2}+\sigma_*^2},\qquad
\xi_0=\Lambda\zeta^*,\qquad
\varphi^*=\exp\Bigl(\Lambda\int_0^\eta\zeta^*(w)\,dw\Bigr).
\end{equation}
Parameter order \lnos{9778--9779}:
\begin{equation}
\label{eq:B40}
\begin{gathered}
M_d,T_d,P_*,\lambda,h
\;\to\;
\text{matching tolerance},\; j_0
\;\to\;
\delta_*,\sigma_*,\Lambda\\
\;\to\;
(B_k),T_{\mathrm{sh}}
\;\to\;
C,X_R
\;\to\;
\kappa_0,t_1,\text{final widths}.
\end{gathered}
\end{equation}

\ORIG
Endpoint test of Prop.~B.3 of [S] \lnos{9281--9285,9353}: at $X_0=4/\Lambda$ there is a fixed
$c_{\mathrm{ex}}>0$ (the proof takes $c_{\mathrm{ex}}=.2$) such that
\begin{equation}
\label{eq:B19}
p_1+\frac{p_2^2}{p_1}>2+c_{\mathrm{ex}}.
\end{equation}
The large-$C$ mechanism requires $\chi\le .99$ and $|Z^*|>\delta_*$, hence $|n_s|\ge\delta_*/2$,
\begin{equation}
\label{eq:p2C}
|p_2|=C\sqrt{\frac{2}{\Lambda}}\,\frac{|n_s|}{\varphi^*\Phi}
\qquad\text{\lnos{9344--9349}}.
\end{equation}
On the complementary region $|Z^*|>\delta_*$ one has $|n_s|\ge\delta_*/2$, so this term carries $C$;
that region does not contain $Z^*(0)=0$\lnos{9344--9352}.

\subsection{Parity: \texorpdfstring{$p_2(\cdot,0)\equiv 0$}{p\_2(.,0)=0} exactly}
\label{sec:note3-parity}

\begin{proposition}[On the axis at $j_0=0$, $p_2\equiv 0$]
\label{prop:p2zero}
\DER
At $j_0=0$, $U^*$ is odd and $\Pi_0$ is even\lnos{8998--9004}
$\Rightarrow$ $Z^*$ is odd $\Rightarrow$ $Z^*(0)=0$;
$\chi$ is even, $\varphi^*$ is even (\eqref{eq:B3}).
The comparison profile of Prop.~B.2 of [S] $\Phi_0=f_0(Y\chi)$ is even,
and $u_0=-YZ^*/(2L)$ is odd\lnos{9236}.
The remainders $R_1$ is even and $R_2$ is odd (substitute each term of \lnos{9197--9199}).
The fixed-point map \lnos{9237--9243} preserves this parity.
The real fixed point is unique\lnos{9261}
$\Rightarrow$ the exact solution satisfies $u(Y,0)=0$, hence $u_Y(Y,0)=0$.
\ORIG
$n_s=-2u_Y$\lnos{9324}, so $n_s(Y,0)=0$,
hence $p_2(4/\Lambda,0)=0$. Thus $p_2^2/p_1=0$.
\end{proposition}

\begin{proof}[Argument sketch]
Check the parity of each term of \lnos{9197--9199} under $\eta\mapsto-\eta$;
the map \lnos{9237--9243} preserves the class; uniqueness\lnos{9261} forces the exact solution into
that class.
\ORIG ``Real data and uniqueness of the fixed point give a real solution'' is the statement of
\lnos{9261}; this note quotes it as is and does not isolate it as a separate lemma; the full term-by-term
parity verification is not written out.
\notproved
\end{proof}

\begin{remark}[A forbidden generic claim]
\DER
One must not write $p_2^2/p_1\asymp C^2\Lambda^{-1}$ at $\eta=0$:
that asymptotics applies only on the complementary region $|n_s|\ge\delta_*/2$\lnos{9344--9352},
which does not contain $Z^*(0)=0$. The missing step is the jump from \eqref{eq:p2C} to
$|n_s|=\Theta(1)$.
Even without parity, citing only (B.13)\lnos{9172--9182} via (B.18)\lnos{9278},
at $\eta=0$ one has $n_s=O(\Lambda^{-1})$ and $p_1=O(\Lambda^{-1})$
(\lnos{9325--9326}: the upper bound for $p_1$ and the approximation for $n_s$ follow from these), so
$p_2\asymp C\Lambda^{-3/2}$, $p_2^2/p_1\asymp C^2\Lambda^{-2}$, still not $\Lambda^{-1}$.
\end{remark}

\subsection{\texorpdfstring{$p_1=2(3-h)/\Lambda$}{p\_1=2(3-h)/Lambda} and the degeneration of (B.19)}
\label{sec:note3-p1}

\begin{proposition}[On-axis $R_1^{(0)}$ and $\partial_Y\Phi_1$]
\label{prop:Phi1}
\DER
At $j_0=0$: $\chi(0)=H^*(0)=\zeta^*(0)=Z^*(0)=0$, $u_0(Y,0)=0$,
$\Phi_0(Y,0)=f_0(0)=1$\lnos{9150--9160}.
$L(0)=1$, $W^*(0)=-3$\lnos{9005,9008}, so
\begin{equation}
\label{eq:R10}
R_1^{(0)}(Y,0)=W^*(0)+h=-3+h,
\end{equation}
independent of $\sigma_*,\delta_*,Y$.
At $\eta=0$, $T=J_2\chi/2=0$, so
$2\bigl(Y(\Phi_1)_{YY}+2(\Phi_1)_Y\bigr)=h-3$.
Smoothness at $Y=0$ excludes the $Y^{-2}$ homogeneous solution, giving
\begin{equation}
\label{eq:Phi1Y}
\partial_Y\Phi_1(4,0)=\frac{h-3}{4}<0
\qquad\bigl(0<h\le 10^{-2},\ \text{\lnos{9002,1344--1347}}\bigr).
\end{equation}
The remainder $O(\Lambda^{-2})$ does not change the sign for $\Lambda$ large.
\end{proposition}

\begin{proof}[Argument sketch]
The fixed-point expansion $\Phi=\Phi_0+\Lambda^{-1}\Phi_1+O(\Lambda^{-2})$ comes from the same
contraction map \lnos{9237--9247}, not from an external fit.
Cross-check: on the axis $S_q=-W-h=3-h$\lnos{1464}, consistent with $R_1=h-3$.
ODE: $Yv'+2v=(h-3)/2$ with $v(0)$ finite $\Rightarrow$ $v\equiv(h-3)/4$.
The consistency estimate of the expansion remainder is not written out.
\notproved
\end{proof}

\begin{proposition}[$p_1(4/\Lambda,0)$]
\label{prop:p1}
\ORIG
$p_1=a=-2Y\Phi_Y/\Phi$\lnos{9324}.
\DER
$\Phi_Y(Y,0)=\Lambda^{-1}(h-3)/4+O(\Lambda^{-2})$,
$\Phi(Y,0)=1+O(\Lambda^{-1})$, so
\begin{equation}
\label{eq:p1}
p_1\Bigl(\frac{4}{\Lambda},0\Bigr)
=\frac{2(3-h)}{\Lambda}+O(\Lambda^{-2})>0
\qquad(\Lambda\text{ sufficiently large}).
\end{equation}
Cross-check: on the axis $Q_s(0)=S_q(0)/2$\lnos{1487}, $S_q(0)\approx 3-h$,
$p_1=XQ_s/L\approx(4/\Lambda)(3-h)/2=2(3-h)/\Lambda$.
\end{proposition}

\begin{proposition}[(B.19) degenerates at $\eta=0$]
\label{prop:B19deg}
\DER
By Prop.~\ref{prop:p2zero}, $p_2(\cdot,0)\equiv 0$, so $C$ does not enter \eqref{eq:B19}.
\eqref{eq:B19} degenerates to
\begin{equation}
\label{eq:B19deg}
p_1>2+c_{\mathrm{ex}}=2.2.
\end{equation}
By \eqref{eq:p1}, $p_1=2(3-h)/\Lambda\not>2.2$ ($\Lambda$ large).
The large $C$ multiplies $p_2$\lnos{9346--9349} and is ineffective against $p_2=0$.
\end{proposition}

\begin{claim}[The two cones are separate]
\label{cl:two-cones}
\DER
First cone inequality $a-b_s w>0$:
$b_s=2D_X U/E$\lnos{1494--1496} vanishes at $\eta=0$ ($U$ odd), $a=p_1>0$,
\emph{so this item passes}.
What fails is the second cone test $v_s=p_1+p_2^2/p_1>2$.
The two must not be conflated into one ``cone''.
The literal conditions $p_1>0$ / $\partial_Y\Phi_1<0$ pass, and hold for all $\sigma_*,\delta_*$;
they are not a necessary and sufficient condition for large $C$ to rescue \eqref{eq:B19}
(one still needs $n_s\neq 0$).
\end{claim}

\begin{remark}[The quantifiers of (B.2)]
\ORIG
The quantifiers of (B.2): first choose $\delta_*>0$ so that $\{|Z^*|\le\delta_*\}$ avoids a
neighborhood of $\eta_0$, then choose $\sigma_*$ so that $\chi>.99$\lnos{9014--9019}.
\DER
Substituting $j_0=0$ (an external limit point; the source allows only $0<j_0\le.05$\lnos{9002}):
$\eta_0=0$ and $Z^*(0)=0$ (\lnos{9004--9008} at $\eta=0$),
so for every $\delta_*>0$ the point $\eta=0$ lies both in $\{|Z^*|\le\delta_*\}$ and in any
neighborhood of $\eta_0$.
\eqref{eq:B2} has no solution under the quantifiers of this dichotomy. This is not a choice problem
repairable by making $\sigma_*,\delta_*$ smaller/larger.
This statement stays at level L0 (this dichotomy fails at the external limit point $j_0=0$,
see Claim~\ref{cl:L0}) and is not extended to an existence statement under other dichotomies.
\end{remark}

\subsection{The L0 conclusion and the in-class rigidity}
\label{sec:note3-L0}

\begin{claim}[L0 (the device extended to the external limit point $j_0=0$)]
\label{cl:L0}
\DER
$j_0=0$ lies outside the source's domain (the source allows only $0<j_0\le.05$\lnos{9002}; this note
extends the device to that external limit point): the unique real fixed point lies on the
$\eta\mapsto-\eta$ symmetry, so $p_2(\cdot,0)\equiv 0$ and $b_s(\cdot,0)=0$;
\eqref{eq:B19} degenerates at $\eta=0$ to $p_1>2.2$, while $p_1=2(3-h)/\Lambda\to 0$.
Hence the endpoint test \eqref{eq:B19} of Prop.~B.3 of [S] fails at that extended point.
This is L0: ``the device fails when extended to $j_0=0$'';
the nonexistence of a $j_0=0$ variant within one class (Prop.~\ref{prop:L1}) is L1,
and this note does not claim that the source itself fails at some step with $j_0>0$.
\end{claim}

\paragraph{The necessity at $X_0$.}
\ORIG
Lemma~B.4 of [S] requires $v_r>2+c$ on the whole interval $[X_0,X_i]$\lnos{9383--9389};
at $\eta=0$, $v_r=p_{1,r}$.
The analytic axis region extends only to $Y\le 4.1$, i.e.\ $X\le 4.1/\Lambda\to 0$\lnos{9166--9167,9366}.
The first collar, where $v_s>2$, is produced by continuity from $v_r(0)>2+c$\lnos{9503--9506}.
\DER
The following proposition records ruling (i): L1 (no $j_0=0$ variant within the class), not L2.

\begin{proposition}[In-class rigidity (L1)]
\label{prop:L1}
Assume the axis device satisfies the following class definition (H1)--(H3):
\begin{itemize}
\item[(H1)] The axis-data class of Appendix~B of [S] ((B.1)--(B.3), \lnos{9003--9036},
$U^*=4\eta+j_0$), the analytic axis profile defined on $0\le Y=\Lambda X\le 4.1$\lnos{9166--9167},
with vanishing leading residual stress on the axis region ($T_0=0$, \lnos{1904--1905,2111}).
\item[(H2)] Stress activated at $X_0=4/\Lambda$: the statement of Prop.~B.5 of [S] (\lnos{9462--9464},
continued to $X_i$ with $X\le X_0$ unchanged), (B.26) (\lnos{9476--9479}, the prescription on
$0<y<t_1$), and $B_0(0,\eta)\ne 0$ (\lnos{9523--9524}).
\item[(H3)] The same cone test: \eqref{eq:B19} of Prop.~B.3 of [S] (\lnos{9281--9285},
$p_1+p_2^2/p_1>2+c_{\mathrm{ex}}$) and the $v_r>2+c$ of Lemma~B.4 of [S] on
$[X_0,X_i]\times[-1,1]$ (\lnos{9384--9389}), with the first collar given by the continuity of
\lnos{9505--9506}, uniformly in $\eta$.
\end{itemize}
\DER
Then $j_0=0$ fails the cone test within this class. Concretely, at $(X,\eta)=(4/\Lambda,0)$:
\begin{enumerate}
\item At $j_0=0$, $U^*$ is odd and $\Pi_0$ is even\lnos{8998--9004} $\Rightarrow$ $Z^*$ odd,
$\chi,\varphi^*$ even, the comparison profile $\Phi_0=f_0(Y\chi)$ even, $u_0=-YZ^*/(2L)$ odd\lnos{9236};
the remainders $R_1$ even and $R_2$ odd\lnos{9197--9199}; the fixed-point map\lnos{9237--9243}
preserves this parity; the real fixed point is unique\lnos{9261} $\Rightarrow$ the exact solution has
$u(Y,0)\equiv 0$
$\Rightarrow$ $n_s=-2u_Y=0$\lnos{9324} $\Rightarrow$ $p_2(4/\Lambda,0)=0$
(\lnos{9277,9346--9349}).
\item $R_1^{\mathrm{lead}}(Y,0)=W^*(0)+h=-3+h$ (\lnos{9005,9008} at
$j_0=0,\eta=0$; the $LR_1$ of (B.14) evaluated at $\eta=0$\lnos{9197}; $L(0)=d(0)=1$), independent of
$\sigma_*,\delta_*,Y$; the axis-regular solution of the slice ODE
$Yv'+2v=(h-3)/2$ gives $\partial_Y\Phi_1(4,0)=(h-3)/4<0$;
hence by Prop.~\ref{prop:p1}, $p_1(4/\Lambda,0)=2(3-h)/\Lambda+O(\Lambda^{-2})\to 0$
(cross-check: the $Q_s(0)=S_q(0)/2$ of \lnos{1487} and the $S_q=-W-h=3-h$ of \lnos{1464}).
\item Hence \eqref{eq:B19} degenerates at $\eta=0$ to $p_1>2.2$ and fails (Prop.~\ref{prop:B19deg}).
The first cone inequality $a-b_sw>0$ passes ($b_s=2D_XU/E$ vanishes at $\eta=0$, \lnos{1494--1496});
what dies is the second cone test $v_s>2$ (Claim~\ref{cl:two-cones}).
\end{enumerate}
\end{proposition}

\paragraph{Evidence for the negative statement (``must start at $X_0$'').}
\ORIG
\begin{enumerate}
\item Prop.~B.5 of [S] \lnos{9499--9500}: for $t_1$ sufficiently small, $P_c>2+c$; by (B.29)
(\lnos{9494--9498}), $P_c=v_r+O(ye_a)$, so at $y=0$, $v_r(X_0,\eta)>2+c$. Shrinking $t_1$ only
shrinks the error term and cannot raise the endpoint where $v_r(0)<2$.
\item Prop.~B.5 of [S] \lnos{9505--9506}: ``Continuity from $v_r(0)>2+c$ gives one positive-width
first collar on which $v_s>2+c/2$, uniformly in $\eta$'' --- this is the only sentence in the whole
source that names $v_r>2+c$ at $y=0$. (B.26) (\lnos{9476--9479}) \emph{does not} use $v_r>2+c$.
\end{enumerate}

\paragraph{Where the chain breaks, and downstream uses.}
\DER
If the endpoint lower bound is deleted, then at $\eta=0$: the $\eta$-uniform $v_s>2$ of the first
collar cannot start (the quantifier of \lnos{9505--9506}) $\to$ the near-endpoint $P_c>2$ fails $\to$
both the admissible and the relaxed branches fail simultaneously (both branches of the dichotomy of
\lnos{9503--9505} break) $\to$ while $T_0$ is already nonzero at $X_0+$ (\lnos{9523--9524}) $\to$
Cor.~B.6 of [S] (\lnos{9563--9564}). This is the load-bearing chain.
Downstream uses (links in the full existence theorem, \emph{not} independent steps of this proposition):
Prop.~4.10(ii) of [S] (\lnos{2112--2129}), Appendix~C of [S] (\lnos{9809--9811}),
Theorem~4.6(iii) of [S] (\lnos{1908--1909}).

\paragraph{Exclusion of candidate gaps.}
\ORIG
\begin{itemize}
\item The $v_s\le 2$ branch of \lnos{9503--9505} still requires $P_c>2$;
\item the $v_s<1$ of \lnos{9538--9544}: its onset is after the first collar, still on the scale
$Y=O(1)$, and the branch follows the reference extension to $X_b$ (and on to $X_i$, \lnos{9366,9540});
\item the vacuum region $T_0=0$ ends exactly at $X_0$ (\lnos{1904--1905,2111});
\item the $4e^{t_c}<4.1$ of Cor.~B.6 of [S] (\lnos{9558}) locks the first collar within the analytic
axis scale, so it cannot be moved out to $X\asymp 1$.
\end{itemize}
\emph{Scope: this proposition is L1 (no $j_0=0$ variant within one function class; $j_0=0$ is an
external limit point outside the source's domain, the source allowing only $0<j_0\le.05$\lnos{9002}),
valid only within the class (H1)--(H3); it is not L2 (it does not concern ``nonexistence of
reflection-symmetric forced blowup''). Outward statements use L0 only.}

\begin{remark}[No assertion outside the class]
\label{rem:outside}
\ORIG
This proposition covers only the class (H1)--(H3), and makes no assertion outside it. Class-external
modifications may change the conclusion, e.g.:
dropping ``the axis profile is unchanged for $X\le X_0$''\lnos{9460--9461} or ``$T_0=0$ on the axis
region''\lnos{1904--1905,2111};
moving the activation point from $X_0$ into the interior of $(X_0,X_i]$ --- Prop.~B.5 of [S] pins only
the strict relaxed cone on the rest of $(X_0,X_i]$\lnos{9461--9464}, and this proposition does not rule
out other activation geometries (the $B_0(0,\eta)\ne0$ of (H2), \lnos{9523--9524}, is precisely the anchor
that pins the activation point);
or changing the cone test. For all such class-external cases this proposition is silent; this is exactly
why it is not written as L2.
\end{remark}

\begin{remark}[Matching window: a single-point check, not a conclusion of this note]
\label{rem:window}
\DER
At the single point $(4/\Lambda,0)$: $H^*(0)=j_0$, $Z^*(0)=-(\tfrac92+h)j_0$ exactly\lnos{9003--9005};
$p_2=-C\sqrt{2/\Lambda}(\tfrac92+h)j_0\,(1+\cdots)$\lnos{9346--9349};
$p_1=2(3-h)/\Lambda+2\chi(0)+\cdots$ (\lnos{9324} together with the $S_q$ and $\chi$ terms
\lnos{1487,9008,9299});
hence $p_2^2/p_1=C^2k^2j_0^2/[(3-h)+\Lambda\chi(0)]$, and when $\chi$ is subdominant the $\Lambda$
cancels.
\HYP
Regime $j_0=o(\sigma_*)$: $\sigma_*$ keeps a positive lower bound independent of $j_0$ --- the source
lines 9014--9015 and 9020--9023 state only ``$\sigma_*$ may be taken sufficiently small'' and ``$\Omega$
exists'', and do not state that the holomorphic-neighborhood width does not collapse as $j_0\to 0$; this
positive lower bound is an additional hypothesis with no source basis.
\HYP
The main case additionally requires $C^2\gg\Lambda/\sigma_*^2$ (the $\chi$ term does not dominate $p_1$)
--- likewise an additional hypothesis with no source basis: the source lines 9014--9015, 9020--9023 state
only that $\sigma_*$ may be taken sufficiently small and that $\Omega$ exists.
\DER
The lower bound for $p_2^2/p_1>2.3$ is $\Theta(1/C)$ (the coarse bound $p_1\le C_1$ loses the
$\Lambda^{-1}$ gain at this point --- that upper bound is true in itself; using it to estimate the ratio
only yields the overly strong $\sqrt{\Lambda}/C$).
The matching bound\lnos{2210--2213}:
$\|G_i-4\eta\|_{C^k}\le C_kj_0+C^{\mathrm{nat}}_k/\Lambda+\cdots$;
requiring this bound to be smaller than $\varepsilon_m$\lnos{2190} gives an upper bound on $j_0$ of order
$\Theta(1/C_k)$.
\ORIG
$C$ has no upper bound: (B.16) (\lnos{9215--9218}) gives only a lower bound; $C$ may continue to
increase (\lnos{9350--9352,9583--9585}).
\DER
Hence the $\Theta(1/C)$ lower bound ``barely bites'': increasing $C$ can push the lower bound
arbitrarily small; the only point where the test truly collapses is the exact value $j_0=0$
(the parity mechanism of Prop.~\ref{prop:L1}).
Ruling (a conclusion of this note): merely ``the matching upper bound and the \emph{single-point} cone
test do not conflict under the parameter order of \eqref{eq:B40} (freeze $j_0$ first, then choose
$\delta_*,\sigma_*,\Lambda,C$)''.
We do not write: a quantitative window theorem; nor ``any $j_0>0$ is rescued by large $C$''
(the $\eta$-neighborhood, the worst point of the (B.2) complementary region, and the whole interval
$[X_0,X_i]$ of Lemma~B.4 of [S] are not covered).
\notproved
\end{remark}

%% =====================================================================
\section{Acknowledgments and statement}
\label{sec:ack}

\paragraph{Acknowledgments.}
The author (Kaiyan Ren) is an M.Sc. student in the Department of Computer Science and Engineering,
School of Engineering, The Hong Kong University of Science and Technology, and was formerly an
undergraduate in the School of Mathematics, Nanjing University.
Gratitude is due to several independent review sessions for their line-by-line scrutiny of the line
anchors, the color-code and the scope boundaries: every \ORIG statement of this note was, on that
account, re-checked verbatim against the extracted text, and several wordings and anchors were corrected.
The review records themselves do not constitute evidence for this note; the only evidence is the source
line numbers.

\paragraph{Statement on the use of artificial-intelligence tools.}
This note was written with the assistance of generative-AI tools: DeepSeek (line-number extraction,
formula rewriting, parity computations and the orchestration of the review rounds), Claude (independent
review and the packaging of the release draft), ChatGPT (earlier finite-dimensional notes and material
preparation). These tools are not authors of this note. All assertions are anchored to line numbers of
the reference text extraction \nolinkurl{navier-stokes.txt} identified in \S\ref{sec:intro}; the reader
should check them line by line, and tool output is not evidence. The author takes full responsibility for the entire content of this note,
including line-number and scope errors.

\paragraph{Scope and disclaimer.}
This note records three structural observations; it is not a completed mathematical paper and has not
been peer-reviewed.
Prop.~\ref{prop:osc}--\ref{prop:V}, \ref{prop:shell}--\ref{prop:rigid},
\ref{prop:p2zero}--\ref{prop:B19deg} give only argument sketches, each explicitly marked as not
constituting a complete proof; steps tagged \HYP are unproven.
The in-class rigidity Prop.~\ref{prop:L1} and the single-point check Rem.~\ref{rem:window} are valid only
within the class definition (H1)--(H3) and the stated single-point check, respectively.
This note passes no judgment on the correctness of the source construction: the whole text reads the
source under the hypothesis ``the construction is correct as claimed'', and the L0 statement concerns
only the extension of the device to $j_0=0$ (an external limit point outside the source's domain).
\textbf{No new Navier--Stokes solution is claimed}, no new regularity criterion is claimed, the existence
or nonexistence of reflection-symmetric forced blowup is not claimed, and the finite-dimensional
five-moment observation or the shell representation of the terminal force $F_0$ is not promoted to a
uniqueness theorem for the full construction.

%% =====================================================================
\begin{thebibliography}{9}
\bibitem{source}
OpenAI,
\textit{Finite Time Blowup for Navier--Stokes}
(these words are the title line of the extraction, \lnos{2--3}; abstract at \lnos{4--6},
Theorem~1.1 of the source at \lnos{28--37}),
made public September 8, 2026.
Announcement and full-text entry: \mbox{\url{https://openai.com/index/navier-stokes-solution/}}.
(As of the finalization of this note the paper is not on arXiv, so it is cited via its official
publication page, without an arXiv number.)
All line numbers of this note refer to the reference text extraction \texttt{navier-stokes.txt} of that
PDF (10264 lines; fingerprint and availability in the reproducibility paragraph of \S\ref{sec:intro}).

\bibitem{lean}
OpenAI,
\textit{Lean certificates accompanying Navier--Stokes and Euler results},
GitHub repository \mbox{\url{https://github.com/openai/NavierStokesAndEuler}},
fixed commit \texttt{8937a8f4cbc7abaab5e9e97d1cc7f5d2319d9538}
(this note does not cite the Lean proofs).
\end{thebibliography}

\end{document}
